% =====================================================================
% 12c_hypotheses.tex — hypothèses de cadrage
% =====================================================================
\chapter{Hypothèses de cadrage}
\label{ann:hypotheses}

\section{Statut des hypothèses}

Le cadrage a identifié quatorze ambiguïtés ou contradictions apparentes dans le
sujet. Chaque hypothèse a été confrontée au comportement réel d'Oracle 21c XE,
puis retenue ou adaptée. Les hypothèses ci-dessous sont historiques : elles
expliquent pourquoi le schéma et les scripts ont pris leur forme actuelle.
Lorsqu'une hypothèse n'a pas été implantée exactement comme prévu, le statut le
précise.

\begin{description}[leftmargin=2.2em,style=nextline]

\item[H1 — Parrainage dès la partie A.]
Le sujet demande deux clients parrainés en A4, mais l'attribut
\code{ref\_parrain} n'apparaît qu'en E10. Il est impossible d'insérer un parrain
avant que l'attribut existe.
\\ \emph{Retenu.} L'\code{ALTER TYPE ... CASCADE} est exécuté en partie A, juste
après la création de \code{client\_type} et avant les insertions. E10 rejoue
l'ajout sous garde et démontre la double notation. Le script reste rejouable.

\item[H2 — CC99 sans titulaire.]
B5b demande d'insérer un compte courant sans titulaire, alors que A3 impose
\code{CHECK (ref\_client IS NOT NULL)}.
\\ \emph{Retenu.} La contrainte \code{ck\_cc\_client} est
\code{DEFERRABLE INITIALLY DEFERRED} sur \code{cpt\_courant}. Elle reste
\code{NOT NULL}, mais la vérification a lieu au \code{COMMIT}. Les contraintes
équivalentes de \code{cpt\_epargne} et \code{cpt\_terme} restent immédiates.

\item[H3 — Plafond des dépôts d'épargne.]
Le sujet évoque un plafond de dépôt sur l'épargne, alors que les opérations
sont réservées aux comptes courants.
\\ \emph{Cadrage.} Un trigger sur \code{cpt\_epargne} avait été envisagé. Dans
le jeu final de six triggers, ce contrôle n'est pas présent : les mouvements
d'épargne ne sont pas créés, et le plafond reste un attribut contrôlé
positivement (\code{plafond\_depot >= 0}). La limite est documentée dans les
perspectives : si les opérations d'épargne étaient ajoutées, il faudrait
vérifier \code{solde + montant <= plafond\_depot}.

\item[H4 — Prélèvement de la commission de virement.]
Le sujet ne précise pas où la commission est prélevée.
\\ \emph{Retenu.} La commission vaut
\code{GREATEST(100, 0,01 * |montant|)} sur l'opération de débit et 0 sur le
crédit. Le solde varie de \code{montant - commission}. La commission est donc
visible sur le compte source.

\item[H5 — Périmètre du plafond journalier.]
Le plafond de 200\,000~DA peut s'interpréter de plusieurs façons.
\\ \emph{Retenu.} Il porte sur la somme des \code{RETRAIT} (montant négatif)
d'un même compte, le même jour (\code{TRUNC(date\_op)}), hors commission. Les
\code{CB} ne sont pas comptés, par lecture littérale du mot « retraits ».

\item[H6 — Formule des intérêts.]
\code{solde * (1 + taux/100)} est le capital capitalisé, pas l'intérêt seul.
\\ \emph{Retenu.} \code{calculeInterets} renvoie ce capital ; l'intérêt affiché
vaut \code{solde * taux/100}. Pour le terme :
\code{solde * (1 + taux/100 * durée\_mois/12)}.

\item[H7 — Sortie anticipée d'un compte d'épargne.]
Un compte d'épargne n'a pas d'échéance.
\\ \emph{Retenu.} Une pénalité de 5~\% est appliquée au solde
(\code{solde * 0.95}) ; sur le compte à terme, elle porte sur le montant
bloqué.

\item[H8 — Compte dormant.]
Aucun trigger ne peut détecter une absence d'opération.
\\ \emph{Retenu.} La procédure \code{maj\_statuts\_dormants} effectue un
balayage ; un job \code{DBMS\_SCHEDULER} optionnel pourrait l'automatiser. La
requête 6.23 fournit la lecture correspondante.

\item[H9 — Intégrité des \code{REF} des signataires.]
\code{ORA-30730} interdit la FK et \code{ORA-22812} interdit le trigger sur la
table de stockage d'une \emph{nested table}.
\\ \emph{Retenu.} L'intégrité est assurée par les procédures d'écriture
(\code{ajout\_client}, \code{drop\_client}, insertion des signataires) et par
une requête de contrôle. Les autres \code{REF} de premier niveau ont une vraie
FK.

\item[H10 — Unicité des numéros de compte.]
Trois tables distinctes imposent trois clés primaires locales.
\\ \emph{Retenu.} Les préfixes \code{CC}, \code{CE} et \code{CT} distinguent
les comptes, avec un contrôle applicatif. L'absence de table mère est
documentée comme limite du modèle objet.

\item[H11 — Requête 6.25 sans effet.]
Avec les seules données de A4, aucun compte à terme n'a simultanément une
échéance dépassée de plus d'un an et un montant bloqué nul.
\\ \emph{Retenu.} Un quatrième compte \code{CT4} est ajouté au jeu de test, avec
une échéance de $-400$ jours et un montant bloqué nul. La requête 6.25 supprime
alors exactement une ligne.

\item[H12 — Alertes d'échéance.]
Le sujet demande une alerte huit jours avant, puis à nouveau si l'échéance est
dépassée.
\\ \emph{Retenu.} Deux types : \code{ECHEANCE\_PROCHE} si
$0 \leq \text{échéance} - \text{SYSDATE} \leq 8$, et
\code{ECHEANCE\_DEPASSEE} si l'échéance est passée. L'insertion est idempotente,
et \code{scan\_echeances} traite les comptes déjà en base.

\item[H13 — Opérations d'un client supprimé.]
La FK \code{operation.ref\_client} empêcherait la suppression d'un client qui a
opéré sur les comptes d'autrui.
\\ \emph{Retenu.} \code{drop\_client} supprime aussi les opérations que ce
client a faites ailleurs. La traçabilité est sacrifiée au profit de la
cohérence référentielle ; cette décision est documentée.

\item[H14 — Mot de passe et création de l'utilisateur.]
\code{MASTER1} existe déjà avant le projet.
\\ \emph{Retenu.} Le script \code{00\_creation\_utilisateur.sql} est fourni mais
n'a pas été exécuté. Il est idempotent. Les scripts supposent
\code{MASTER1/system} ; le mot de passe réel doit être substitué lors du
rendu.

\end{description}

\begin{remarquebox}[Lecture des hypothèses]
Une hypothèse n'est pas une erreur : c'est une décision prise face à une
ambiguïté. En soutenance, il faut pouvoir dire laquelle a été retenue, pourquoi,
et quelle alternative a été écartée. Les hypothèses H1, H2 et H11 sont les plus
structurantes : elles expliquent l'\code{ALTER TYPE} en partie A, la contrainte
différée de B5b et le quatrième compte à terme.
\end{remarquebox}
